\section{Conclusion}
\label{sec:conclusion}

\subsection{The Container Parallel}

Malcolm McLean did not invent cargo boxes. Wooden crates, barrels, and pallets had been moving
cargo for centuries before April 26, 1956. What McLean invented was the \textit{interface
standard}: the specification that made any container interchangeable with any vessel, any crane,
and any chassis, without negotiation, without custom handling, and without carrier-specific
equipment. The value was not in the box; it was in the standardization. Once every ship used
the same corner fittings, the cargo logistics industry reorganized around the interface, and
global shipping costs fell 40--60\% over the subsequent three decades \citep{Levinson2006}.

The relay marketplace is McLean's insight applied to the American trucking industry. Relay
driving is not new. The value is not in the concept; UPS, BNSF, and FedEx Ground have been
executing relay operations internally for decades. What does not exist is the interface standard
that makes any relay driver interchangeable with any truck, any load, and any carrier --- the
RDU specification, the standardized ELD handoff API, the hub slot system, and the insurance
framework that travels with the driver rather than the carrier. Without these standards,
every relay arrangement is bespoke, bilateral, and costly. With them, any carrier can relay
any load through any hub using any available relay driver from the pooled workforce, at a
known cost per swap, with known insurance coverage, on a known schedule.

The 65-year gap between the interstate system's completion and the relay marketplace is not
evidence that relay is impractical. It is evidence that coordination problems do not solve
themselves. McLean's containerization required shipping companies to agree on corner fittings
before any individual ship could recoup the investment in container cranes. The relay
marketplace requires FMCSA to define three regulatory terms before any individual hub operator
can recoup the investment in relay infrastructure. The parallel is exact.

\subsection{What the Rulemaking Needs to Say}

The regulatory work required to launch the relay marketplace is modest in scope, precedented
in form, and achievable within the standard FMCSA rulemaking timeline of 18--24 months.
Three regulatory definitions are required:

\begin{enumerate}
  \item \textbf{Relay terminal (49 CFR Part 395, new \S 395.XX)}: A ``relay terminal'' is a
  FMCSA-designated facility at which a motor carrier may transfer a commercial vehicle from
  one qualified CDL driver to another without initiating a new trip for purposes of Hours of
  Service recordkeeping. A relay terminal must be: (a) physically fixed, with dock facilities
  and driver rest amenities; (b) operated by a licensed relay terminal operator holding a
  commercial auto umbrella policy of not less than \$5 million; (c) equipped with an FMCSA-
  certified relay ELD handoff system compliant with the relay API specification. Relay
  terminals are approved by FMCSA upon application; approval constitutes NFRZ status under
  IIJA \S 167 \citep{FMCSA_HOS2022}.

  \item \textbf{Relay pre-trip inspection (49 CFR Part 396, amendment)}: At an approved relay
  terminal, a driver receiving a commercial vehicle from a departing driver may perform a
  condensed relay pre-trip inspection in lieu of the standard 396.11 pre-trip, provided that:
  (a) the vehicle's OBD system returns a green status (no active fault codes, brake pressure
  $\geq$ specification, tire inflation within 10\% of specification, lighting circuit intact) on
  the relay terminal's diagnostic interface; (b) the incoming driver physically verifies the
  5-item relay checklist (coupling security, trailer seal, load securement, fifth-wheel latch,
  emergency brake). A green OBD status plus 5-item physical verification constitutes a complete
  pre-trip inspection for Hours of Service and liability purposes.

  \item \textbf{Relay engagement period (49 CFR Part 395, new \S 395.YY, and 49 CFR Part 387
  amendment for Mode 2)}: A ``relay engagement period'' begins when a licensed relay platform
  confirms an IC relay driver's assignment to a specific relay slot and ends when the driver
  completes the relay handoff at the destination terminal. During the relay engagement period,
  the relay platform's gap insurance coverage applies as primary for IC driver occupational
  injury and secondary for third-party auto liability (with the load carrier's commercial auto
  policy as primary). The relay engagement period definition is the regulatory analog to the
  Transportation Network Company ``Period 2/3'' coverage framework established under state TNC
  insurance laws.
\end{enumerate}

With these three definitions in place, the relay marketplace can launch commercially. No new
statute is required. No highway construction is required. No new enforcement mechanism is
required --- the FMCSA's existing ELD audit infrastructure and carrier compliance program
covers relay terminals under the same framework it applies to any motor carrier operation.
The regulatory investment is three definitions in two existing CFR parts. The commercial
investment is \$5 million per hub, with payback in days.

\subsection{What Builds First}

The sequencing of relay marketplace development follows a four-step initialization order,
from least to most capital-intensive:

\textbf{Step 1: FMCSA Rulemaking NPRM (2025--2026, \$0 capital}.
FMCSA issues an NPRM defining relay terminal, relay pre-trip, and relay engagement period.
The comment period surfaces carrier, insurer, and driver organization input. The rule is
finalized in 2027.

\textbf{Step 2: Technology Stack (2025--2026, \$5--15 million}.
The relay platform technology --- slot exchange, ELD handoff API, driver verification system,
insurance clearance, settlement --- is built by a private technology company (or carrier
consortium) in advance of the rulemaking's finalization. The technology does not require the
rule to be operational; it requires the rule to be commercially deployed at scale. A beta
network of 3--5 hubs on one corridor (NY--LA, starting with Chicago and Denver as the first
two hubs) validates the technology before national rollout.

\textbf{Step 3: Pilot Hub Network (2026--2028, \$25--40 million}.
Five hubs on the NY--LA corridor --- approximately every 400 miles: Chicago (I-80/I-90 junction),
Des Moines (I-80), North Platte (I-80), Salt Lake City (I-80/I-15), and Reno (I-80) ---
provide a 2,800-mile relay-capable corridor. These hubs operate under Mode 1 (W-2 relay
drivers) using the beta regulatory framework. Each hub costs \$5 million; total corridor
investment is \$25 million. Hub revenue from day one provides operating capital.

\textbf{Step 4: National Network Expansion (2028--2032, \$200--400 million}.
Following regulatory finalization and pilot hub validation, the relay marketplace operator
(or a hub operator franchise network) expands to all T1 diamond hub locations designated by
the ROUTE tier analysis \citep{ROUTE_E1}: Atlanta (I-75/I-85/I-20), Dallas (I-35/I-45/I-20),
Kansas City (I-70/I-35), Memphis (I-40/I-55/I-240), and additional T1 nodes. Mode 2 IC gig
layer activates following FMCSA Mode 2 rulemaking. The national relay marketplace is
operational by 2030 on the primary T1 corridors, with expansion to T2 corridors by 2032.

\subsection{Relay as Layer 0 of Interstate 2.0}

The Interstate 2.0 framework \citep{ROUTE_E2} describes a \$253 billion infrastructure
upgrade program --- managed freight lanes, T1 diamond hubs, Donner Pass hardening, and AV
design standards --- as the long-run investment case for American freight competitiveness.
The relay marketplace is Layer 0 of that program: the institutional infrastructure that must
exist before physical infrastructure can fully function.

A T1 diamond hub without a relay marketplace is a rest area. With a relay marketplace, it is
the node of a continuous-operation freight network. Managed freight lanes without relay driver
pools at the hub nodes deliver PTI improvement but not the 48-hour SLA that requires near-
continuous truck operation. AV trucking without relay hubs (as AV handoff nodes) requires
a parallel infrastructure investment in AV-specific terminals that duplicates what relay hubs
already provide.

The ROUTE simulation \citep{ROUTE_C1} is unambiguous: relay alone captures 90\% of the
freight reliability improvement that the full \$121 billion managed-lane program delivers,
at 0.02\% of the cost. This is not an argument against managed lanes --- it is an argument for
sequencing. Relay in 2027. Managed lanes in 2032. AV deployment in 2035. The relay marketplace
buys the 15 years of lead time that highway construction requires. It is not a substitute for
Interstate 2.0; it is the down payment.

\subsection{Summary of Policy Recommendations}

\begin{enumerate}
  \item \textbf{FMCSA}: Issue NPRM defining relay terminal, relay pre-trip, and relay
  engagement period under 49 CFR Parts 395, 396, and 387. Finalize within 24 months.

  \item \textbf{FHWA/USDOT}: Issue NFRZ designation authority under IIJA \S 167, designating
  approved relay terminals as NFRZ facilities with priority NHPP funding eligibility and
  coordination status. The NFRZ designation removes the first-mover investment risk for
  private hub operators.

  \item \textbf{State DOTs}: Coordinate with FMCSA on relay terminal permitting. Relay
  terminals require commercial zoning, stormwater management, and commercial vehicle access
  road permits; state expedited review for NFRZ-designated facilities reduces the permitting
  timeline by 12--18 months.

  \item \textbf{FMCSA (AV)}: In the relay terminal rulemaking, include a future-compatibility
  provision specifying that relay terminal AV upgrade standards (V2I, charging, monitoring
  workstations) are eligible for the same NFRZ coordination status as Mode 1/2 relay
  operations. This avoids a second rulemaking cycle when AV deployment reaches commercial
  scale on T1 managed lanes.

  \item \textbf{Industry}: Tier 1 carriers (J.B. Hunt, Werner, Amazon Freight, Schneider,
  XPO) should form a relay terminal consortium to fund the pilot hub network (\$25 million
  across 5--7 Tier 1 carriers), signal demand to hub operators, and provide operational
  data for the FMCSA rulemaking comment period. The relay marketplace's business case is
  strongest for large carriers whose fleets generate sufficient volume to benefit from
  reserved Tier 1 slots and whose technology organizations can integrate the relay platform
  API at scale.
\end{enumerate}

The relay marketplace is not a moonshot. It is a commercial platform for a service that the
market would have built decades ago if the regulatory framework and the first-mover coordination
mechanism had existed. McLean did not wait for governments to mandate containerization; he
built a container vessel, proved the economics, and watched the industry reorganize around
the interface standard. The relay marketplace starts the same way: one corridor, five hubs,
a working platform. The rest follows.
